
\subsection{Second-kind period lattices and the CM quotient}

\begin{result}{A.4.34}{The exact quotient lattice and its CM case}{Proved}
Let $E/\Qb$, $\omega_i,\eta_i,h,c,A$ have the conventions of A.4.1,
and put
\[
 \mu_i(a)=\eta_i-a\omega_i,\qquad
 \Lambda_a=\Z\mu_1(a)+\Z\mu_2(a).
\]
For $a\in\C$, this is a full lattice in $\C$ exactly when
$|h-a|\ne c$. Off that circle its basis ratio satisfies
\[
 r_a=\frac{\eta_2-a\omega_2}{\eta_1-a\omega_1},
 \qquad
 \operatorname{Im}r_a
 =\frac{A\bigl(|h-a|^2-c^2\bigr)}
        {|\eta_1-a\omega_1|^2}.
\]
There is a holomorphic map
\[
 F_a:E(\C)\setminus\{O\}\longrightarrow \C/\Lambda_a,
 \qquad [u]\longmapsto[\zeta(u)-au],
 \qquad F_a^*(dz)=-(x+a)\omega.
\]
It has no holomorphic extension at $O$ and is not algebraic.
For algebraic $a$ and a full lattice $\Lambda_a$, the complex elliptic
curve $\C/\Lambda_a$ has CM if and only if $E$ has CM and $a=h$.
\end{result}

\begin{proof}
The period formula of A.4.1 gives the exact real-linear map
\[
 T_a:\C\longrightarrow\C,\qquad
 z\longmapsto(h-a)z+c\bar z,
 \qquad T_a(\Lambda)=\Lambda_a.
\]
Writing $h-a=b_1+ib_2$ in real coordinates gives its matrix
\[
 \begin{pmatrix}b_1+c&-b_2\\ b_2&b_1-c\end{pmatrix},
 \qquad \det_{\R}T_a=|h-a|^2-c^2.
\]
Off the stated circle this is an invertible real-linear map, so it
carries the full lattice $\Lambda$ to a full lattice. On the circle it
has rank one, because $c>0$, and its image cannot be a full lattice in
$\C$. No claim of discreteness is made for that rank-one-line image.
Off the circle $\mu_1(a)\ne0$. The oriented-area identity
\[
 \operatorname{Im}\bigl(\overline{\mu_1(a)}\mu_2(a)\bigr)
       =A\det_{\R}T_a
\]
then proves the displayed formula for $\operatorname{Im}r_a$.
In particular,
\[
 \Lambda_h=c\bar\Lambda,\qquad r_h=\bar\tau,\qquad
 j(\Lambda_h)=\overline{j(E)}.
\]
The sign of $\operatorname{Im}r_h$ is negative; reversing the
orientation of the chosen lattice basis gives the usual upper
half-plane convention for its $j$-invariant.

The quasi-period law
$\zeta(u+\lambda)-\zeta(u)=\eta_\lambda$ proves that $F_a$ descends to
$E(\C)\setminus\{O\}$. Since $\zeta'=-\wp$ and $\wp=x$, it also gives
$F_a^*(dz)=-(x+a)\omega$. In a local parameter at $O$, this differential
has a double pole with nonzero leading coefficient. A holomorphic
extension of $F_a$ would pull the holomorphic differential $dz$ on the
target elliptic curve back to a holomorphic differential, which is
impossible. An algebraic map from the punctured smooth curve to this
proper elliptic curve would extend at $O$ by the valuative criterion
of properness. Therefore $F_a$ is not algebraic.

Now assume $a\in\Qb$. If $E$ is non-CM, its four numbers
$\omega_1,\omega_2,\eta_1,\eta_2$ are linearly independent over $\Qb$.
Here is the precise period-theorem input. Apply
\cite[Lemma 15.6, pp.136--138 of the opened manuscript]{EllHuberWustholz2022}
to the $1$-motive $[0\to E]$. In Definition~15.1 it is reduced and
saturated: both auxiliary lattice and torus are zero, and its
endomorphism algebra is $\operatorname{End}^0(E)=\Q$.
Its Betti basis therefore has two elements and its de Rham basis has
two elements. The toric and lattice rows in the lemma are absent;
the lemma consequently asserts independence of the four entries of
the resulting $2$ by $2$ period matrix. Changing its second row from
$(-\eta_1,-\eta_2)$ to $(\eta_1,\eta_2)$ preserves that independence.
If the quotient torus had CM, its lattice ratio $r_a$ would be
imaginary quadratic. The relation
\[
 \eta_2-r_a\eta_1-a\omega_2+a r_a\omega_1=0
\]
would then be a nonzero linear relation over $\Qb$, a contradiction.
This proves that every such quotient is non-CM when $E$ is non-CM.

Suppose instead that $E$ has CM. Then $\tau$ is imaginary quadratic
and $h$ is algebraic by A.4.2. Put
\[
 U=\omega_1,\qquad
 V=c\bar\omega_1=\frac{\pi}{U\operatorname{Im}\tau},\qquad
 b=h-a\in\Qb.
\]
The two numbers $U,V$ are linearly independent over $\Qb$.
Indeed, a dependence with both coefficients nonzero would imply
$\pi=\alpha U^2$ for $\alpha\in\Qb$, contrary to the algebraic
independence of $U$ and $\pi$ in A.4.3, proved using
\cite[Corollary 2, p.343]{EllChudnovsky1978}.
A dependence with one coefficient zero is excluded by $U,V\ne0$.
The quasi-period identity now becomes
\[
 \mu_1(a)=bU+V,\qquad
 \mu_2(a)=b\tau U+\bar\tau V.
\]
If the quotient is CM, $r_a$ is algebraic, and hence
\[
 b(\tau-r_a)U+(\bar\tau-r_a)V=0.
\]
Independence gives $r_a=\bar\tau$ and
$b(\tau-\bar\tau)=0$. Since $\tau\ne\bar\tau$, this forces
$a=h$. Conversely, $a=h$ gives $\Lambda_h=c\bar\Lambda$;
conjugation preserves the nontrivial imaginary-quadratic
endomorphism algebra, so this quotient has CM.
\end{proof}

\paragraph{The remaining quotient arithmetic.}
The proved classification concerns the CM property of the quotient.
For a non-CM $E$, its actual quotient at $a=h$ has the algebraic
$j$-invariant $\overline{j(E)}$ and is itself non-CM. The displayed
maps and period-theorem calculation do not exclude the algebraicity
of this particular $a$. The H1 gate therefore remains unchanged.
